\textbf{Specifiche statiche:} \\
Segnali a gradino sono spesso usati per generare riferimenti
costanti. Per raggiungere l'inseguimento perfetto di un riferimento costante $r(t)$ per la variazione dell'uscita $y(t)$ è necessario imporre:
\[\lim_{t\to\infty} e(t)= \lim_{t\to\infty} (r(t)-y(t))=0 \]
Tuttavia tale condizione è verificata se $G(s)$ ha almeno un polo nell'origine, come nel nostro caso. \\ \\
\textbf{Specifiche dinamiche:} \\
Dobbiamo assicurare un inseguimento perfetto di un gradino con un tempo di risposta non superiore a 45s ed eventuali oscillazioni entro il 5\% dal valore di regime.
\\ Per farlo calcoliamo la $G_c(s)$. \\
Iniziamo utilizzando un'approssimazione ad un polo dominante ovvero: $$G_c(s)\approx \frac{1}{1+s\tau}$$
\begin{comment}
Ricaviamo lo smorzamento $\delta$ sfruttando la seguente relazione: 
\[S=e^{-\frac{\pi\delta}{\sqrt{1-\delta^2}}}\]
Svolgendo i seguenti passaggi: \\
\begin{displaymath}
\begin{array}{c}
\ln(S)=-\frac{\pi\delta}{\sqrt{1-\delta^2}} \\
(1-\delta^2)\ln^2(S)=\pi^2\delta^2 \\
\ln^2(S)-\delta^2\ln^2(S)=\pi^2\delta^2 \\
\ln^2(S)=\delta^2(\pi^2+\ln^2(S)) \\
\delta^2=\frac{\ln^2(S)}{\pi^2+\ln^2(S)} \\
\delta=\sqrt{\frac{\ln^2(S)}{\pi^2+\ln^2(S)}} \\
\end{array}
\end{displaymath}
E sostituendo i valori numerici otteniamo: \\ \\
$\delta=\sqrt{\frac{\ln^2(0.05)}{\pi^2+\ln^2(0.05)}}= 0.69010 \approx 0.7$ \\
\end{comment}
E' possibile ricavare il margine di fase $M_f$ richiesto sapendo che per i sistemi a un polo dominante la $\phi(t)=-90°$, ottenendo quindi $M_f=\pi+\phi (t) = 90°$.
\\
Ricaviamo adesso la pulsazione di taglio $\omega_T$ tramite la seguente relazione: \\ $$\omega_T \geq \frac{3}{T_a} = \frac{3}{45} = 0.0667 \, rad/s$$ \\
Si dovrà quindi ottenere un margine di fase $M_f \approx 90°$ e una frequenza di taglio $\omega_T \geq 0.0667 \, rad/s$.